\documentclass[aspectratio=169]{beamer}
\usetheme{Madrid}
\usecolortheme{default}
\usepackage[utf8]{inputenc}
\usepackage[T1]{fontenc}
\usepackage{hyperref}
\usepackage{booktabs}

% ── Metadata ──────────────────────────────────────────────────────────────────
\title[Short Title]{Full Presentation Title}
\subtitle{Conference / Workshop Name}
\author[J. Smith]{Jane Smith}
\institute[MIT]{Massachusetts Institute of Technology \\ \texttt{jsmith@mit.edu}}
\date{May 2026}

% ── Title slide ───────────────────────────────────────────────────────────────
\begin{document}

\begin{frame}
  \titlepage
\end{frame}

% ── Table of contents ─────────────────────────────────────────────────────────
\begin{frame}{Outline}
  \tableofcontents
\end{frame}

% ── Section 1 ─────────────────────────────────────────────────────────────────
\section{Introduction}

\begin{frame}{Motivation}
  \begin{itemize}
    \item Problem statement: why does this research matter?
    \item Prior work and its limitations.
    \item Our key contributions.
  \end{itemize}
\end{frame}

\begin{frame}{Research Question}
  \begin{block}{Core Question}
    How can we improve X by leveraging Y?
  \end{block}

  \vspace{0.5em}
  \begin{exampleblock}{Hypothesis}
    We hypothesise that Z leads to significant gains in W.
  \end{exampleblock}
\end{frame}

% ── Section 2 ─────────────────────────────────────────────────────────────────
\section{Methodology}

\begin{frame}{Approach Overview}
  \begin{columns}
    \column{0.5\textwidth}
    \textbf{Step 1: Data Collection}
    \begin{itemize}
      \item Dataset A: 50{,}000 samples
      \item Dataset B: 10{,}000 samples
      \item Quality filtering applied
    \end{itemize}

    \column{0.5\textwidth}
    \textbf{Step 2: Model Training}
    \begin{itemize}
      \item Architecture: Transformer-based
      \item Training time: 24 h on 8 GPUs
      \item Hyperparameter search: grid
    \end{itemize}
  \end{columns}
\end{frame}

% ── Section 3 ─────────────────────────────────────────────────────────────────
\section{Results}

\begin{frame}{Quantitative Results}
  \begin{table}
    \centering
    \begin{tabular}{lcc}
      \toprule
      \textbf{Method} & \textbf{Accuracy} & \textbf{F1} \\
      \midrule
      Baseline        & 72.3\%            & 0.68        \\
      Prior SOTA      & 81.5\%            & 0.79        \\
      \textbf{Ours}   & \textbf{87.2\%}   & \textbf{0.85} \\
      \bottomrule
    \end{tabular}
    \caption{Comparison on benchmark dataset.}
  \end{table}
\end{frame}

% ── Section 4 ─────────────────────────────────────────────────────────────────
\section{Conclusion}

\begin{frame}{Summary}
  \begin{itemize}
    \item We proposed a novel approach to X using Y.
    \item Achieved state-of-the-art results on benchmark Z.
    \item Code and data available at \url{https://github.com/username/repo}.
  \end{itemize}

  \vspace{1em}
  \textbf{Future work:} Extend to domain D, improve efficiency.
\end{frame}

\begin{frame}[plain]
  \centering
  {\Huge Thank You!} \\[1em]
  {\large Questions?} \\[2em]
  \href{mailto:jsmith@mit.edu}{jsmith@mit.edu}
\end{frame}

\end{document}
